% 连通分量：三角形加孤立点；单向链1->2->3的强、弱分量；有向环的单一强分量。
\begin{tikzpicture}[>=Stealth,
  v/.style={circle,draw=FigureStroke,fill=FigureBlueFill,line width=.8pt,minimum size=6.5mm,inner sep=0pt},
  e/.style={draw=FigureStroke,line width=.8pt},
  a/.style={draw=FigureStroke,line width=.8pt,->,shorten >=1pt},
  grp/.style={draw=FigureBlue,fill=FigureBlueFill,fill opacity=0,rounded corners=6pt,line width=.7pt},
  sgrp/.style={draw=FigureGreen,dashed,line width=.7pt,rounded corners=6pt},
  note/.style={font=\FigureNoteFont\fontsize{7.5}{9}\selectfont,text=FigureMuted,align=center}]
\begin{scope}
  \node[v] (1) at (0,.7) {$1$}; \node[v] (2) at (-.75,-.5) {$2$}; \node[v] (3) at (.75,-.5) {$3$};
  \node[v] (4) at (1.9,.1) {$4$};
  \draw[e] (1)--(2)--(3)--(1);
  \draw[grp] (-1.25,-.95) rectangle (1.25,1.15);
  \draw[grp] (1.5,-.3) rectangle (2.3,.5);
  \node[align=center] at (.5,-1.55) {(a) 两个连通分量};
\end{scope}
\begin{scope}[xshift=3.8cm]
  \node[v] (1) at (0,0) {$1$}; \node[v] (2) at (1.1,0) {$2$}; \node[v] (3) at (2.2,0) {$3$};
  \draw[a] (1)--(2); \draw[a] (2)--(3);
  \foreach \p in {1,2,3}{\draw[sgrp] ($(\p)+(-.45,-.45)$) rectangle ($(\p)+(.45,.45)$);}
  \draw[grp] (-.65,-.65) rectangle (2.85,.65);
  \node[note] at (1.1,1.0) {虚线：强分量；实线：弱分量};
  \node[align=center] at (1.1,-1.3) {(b) 单向链：三个强分量，\\一个弱分量};
\end{scope}
\begin{scope}[xshift=8.1cm]
  \node[v] (1) at (0,.7) {$1$}; \node[v] (2) at (-.75,-.5) {$2$}; \node[v] (3) at (.75,-.5) {$3$};
  \draw[a] (1)--(2); \draw[a] (2)--(3); \draw[a] (3)--(1);
  \draw[sgrp] (-1.25,-.95) rectangle (1.25,1.15);
  \node[align=center] at (0,-1.55) {(c) 有向环：一个强分量};
\end{scope}
\end{tikzpicture}
